\section{Models of Dynamic Systems}

Model predictive control relies on a mathematical model of the system to predict future behavior. This is one of the fundamental ingredients of the entire framework: without a model, there is no basis for forecasting the effect of future control inputs and therefore no basis for optimization-based planning. In practice, one also needs a state estimator, a suitable optimal control problem formulation, a numerical optimization procedure, and finally a verification that the resulting closed-loop system behaves as desired. In this sense, system theory provides the language and tools on which MPC is built.\\

Dynamic-system models can be classified in several ways. One may distinguish between state-space and transfer-function descriptions, between linear and nonlinear models, between time-varying and time-invariant dynamics, between continuous-time and discrete-time formulations, and between deterministic and stochastic models. In this handout, as in the lecture, the focus is mainly on deterministic models. Physical systems derived from first principles are typically described by nonlinear, time-invariant, continuous-time state-space models, whereas the standard models used in MPC are most often linear, time-invariant, discrete-time state-space models. A substantial part of the modeling workflow therefore consists in transforming the former into the latter in a meaningful way.

\subsection{Nonlinear continuous-time state-space models}

A very general class of dynamical systems can be written in \textbf{continuous-time state-space form} as
\[
\dot{x} = g(x,u), \qquad y = h(x,u),
\]
where $x \in \R^n$ is the state vector, $u \in \R^m$ is the input, and $y \in \R^p$ is the output. The map $g :\mathbb{R}^n \times \mathbb{R}^m \rightarrow \mathbb{R}^n$ describes the system dynamics, while $h :\mathbb{R}^n \times \mathbb{R}^m \rightarrow \mathbb{R}^p$ specifies which quantities are measured or observed. This representation is extremely flexible and can describe a very broad range of physical systems.\\
An important reason why the state-space formalism is so useful is that higher-order ordinary differential equations can always be rewritten as systems of first-order equations. If one starts from an $n$th-order differential equation
$$ x^{(n)} + g_n(x,\dot{x},\ddot{x},\dots,x^{(n-1)}) = 0 $$suitable state variables can be introduced by stacking the variable and its derivatives up to order $n-1$. This is possible by defining
$$x_{j+1} = x^{(j)}, \qquad j=0,1,\dots,n-1,$$
where $x^{(0)}$ is simply $x$. The original equation can then be rewritten as
\[\begin{bmatrix}\dot{x}_1\\ \dot{x}_2\\ \vdots \\ \dot{x}_{n-1}\\ \dot{x}_n\end{bmatrix} = \begin{bmatrix} x_2 \\ x_3 \\ \vdots \\ x_n \\ -g_n(x_1,x_2,\dots,x_n) \end{bmatrix}.\]
The original equation is then transformed into $n$ coupled first-order equations. This conversion is not merely formal: it provides a standard representation in which analysis and controller design can be carried out systematically.\\

The \textbf{drawback} of such nonlinear descriptions is that, although they are general and physically meaningful, their analysis and controller synthesis are often difficult. For this reason, one often replaces them locally by linear approximations around a chosen operating point.

\subsection{Linear time-invariant continuous-time models}

A continuous-time linear time-invariant system is written as
\[
\dot{x} = A_c x + B_c u, \qquad y = Cx + Du,
\]
where $A_c$, $B_c$, $C$, and $D$ are constant matrices of compatible dimensions. Compared with the nonlinear case, this model class is much more \textbf{restrictive}, but it has a decisive \textbf{advantage}: a vast mathematical theory is available for its analysis and control design.\\
One important feature of linear systems is that their solution can be written explicitly. If
\[
\dot{x}(t) = A_c x(t) + B_c u(t), \qquad x(t_0)=x_0,
\]
then the state at time $t$ is given by
\[
x(t) = e^{A_c(t-t_0)}x_0 + \int_{t_0}^{t} e^{A_c(t-\tau)} B_c u(\tau)\, d\tau,
\]
where the matrix exponential is defined by
\[
e^{A_c t} := \sum_{n=0}^{\infty}\frac{(A_c t)^n}{n!}.
\]
This expression makes the effect of the initial condition and the input signal fully explicit.

\subsection{Linearization around an operating point}

Most real systems are nonlinear, yet linear models are much easier to analyze and use in controller design. The standard compromise is therefore to \textbf{linearize the nonlinear dynamics around an operating point}. The idea is that if the system is to operate near a desired equilibrium, then a first-order approximation may be sufficiently accurate in a neighborhood of that point.\\

The mathematical basis is the first-order Taylor expansion. If a nonlinear function is expanded around a point $\bar{x}$, then near $\bar{x}$ it can be approximated by its value at $\bar{x}$ plus the Jacobian times the deviation from that point. The first order Taylor approximation of a function $f(\cdot)$ around $\bar{x}$ is
$$f(x) \approx f(\bar{x}) + \frac{\partial f}{\partial x^\top}\bigg|_{x=\bar{x}} (x-\bar{x} \qquad \text{with} \frac{\partial f}{\partial x^\top} = \begin{bmatrix}
\frac{\partial f_1}{\partial x_1} & \cdots & \frac{\partial f_1}{\partial x_n} \\
\vdots & \ddots & \vdots \\
\frac{\partial f_n}{\partial x_1} & \cdots & \frac{\partial f_n}{\partial x_n}
\end{bmatrix}$$
Applied to a nonlinear state-space system
\[
\dot{x} = g(x,u), \qquad y = h(x,u),
\]
and linearized around a stationary operating point $(x_s,u_s)$ satisfying
\[
g(x_s,u_s)=0,
\]
one obtains a linear model in deviation variables:
$$ \dot{x} = \underbrace{g(x_s,u_s)}_{=0} + \underbrace{\frac{\partial g}{\partial x^\top}\bigg|_{x=x_s,u=u_s}}_{=A_c} \underbrace{(x-x_s)}_{\Delta x} + \underbrace{\frac{\partial g}{\partial u^\top}\bigg|_{x=x_s,u=u_s}}_{=B_c} \underbrace{(u-u_s)}_{\Delta u} $$
Thus
$$\dot{x} - \underbrace{\dot{x}_s}_{=0} = \Delta \dot{x} = A_c \Delta x + B_c \Delta u$$
Proceeding in the same way for the output equation, one obtains
$$\Delta y = \underbrace{\frac{\partial h}{\partial x^\top}\bigg|_{x=x_s,u=u_s}}_{=C} \underbrace{(x-x_s)}_{\Delta x} + \underbrace{\frac{\partial h}{\partial u^\top}\bigg|_{x=x_s,u=u_s}}_{=D} \underbrace{(u-u_s)}_{\Delta u}$$
Thus we obtain the linearized model
\[
\Delta \dot{x} = A_c \Delta x + B_c \Delta u, \qquad \Delta y = C \Delta x + D \Delta u,
\]
where
\[
A_c = \left.\frac{\partial g}{\partial x^\top}\right|_{x=x_s,u=u_s},
\qquad
B_c = \left.\frac{\partial g}{\partial u^\top}\right|_{x=x_s,u=u_s},
\]
\[
C = \left.\frac{\partial h}{\partial x^\top}\right|_{x=x_s,u=u_s},
\qquad
D = \left.\frac{\partial h}{\partial u^\top}\right|_{x=x_s,u=u_s}.
\]
After linearization, it is common to drop the $\Delta$ notation and simply denote deviations again by $x$, $u$, and $y$.\\

The pendulum example clearly illustrates both the usefulness and the limitations of this procedure. Around a chosen equilibrium, such as a constant nonzero angle with zero angular velocity, the nonlinear sine term is replaced by its local linear approximation. The resulting matrices depend on the selected operating point. The linearized dynamics provide a good approximation only in a small neighborhood: for small angles and velocities the behavior of the linear and nonlinear systems is close, whereas farther away the approximation deteriorates. This local nature of linearization is essential to keep in mind in MPC modeling.

\begin{theorembox}[Linearization]
Linearization replaces a nonlinear model by a local linear approximation around a stationary operating point. It is extremely useful for analysis and design, but its validity is inherently local: it can only be trusted in a sufficiently small neighborhood of the linearization point.
\end{theorembox}

\subsection{Discrete-time state-space models}

Since controllers are typically implemented on digital hardware, MPC is naturally formulated in discrete time. A \textbf{discrete-time system} is described by difference equations of the form
\[
x(k+1) = g(x(k),u(k)), \qquad y(k) = h(x(k),u(k)),
\]
where the input and output are defined only at discrete instants indexed by $k \in \mathbb{Z}_+$. Some systems are inherently discrete, such as an account balance updated every month, but most control systems originate in continuous time and are then discretized for implementation.\\
The sampling instants are typically of the form
\[
t_{k+1} = t_k + T_s,
\]
where $T_s$ is the sampling time. A standard assumption is that the control input is held constant between two consecutive sampling instants,
\[
u(t)=u(t_k), \qquad t\in [t_k,t_{k+1}),
\]
which is the usual zero-order-hold implementation.

\subsection{Euler discretization}

For nonlinear systems, one simple discretization method is the forward Euler approximation. Starting from the continuous-time model
\[
\dot{x}_c(t)=g_c(x_c(t),u_c(t)), \qquad y_c(t)=h_c(x_c(t),u_c(t)),
\]
one approximates the derivative by
\[
\dot{x}_c(t) \approx \frac{x_c(t+T_s)-x_c(t)}{T_s}.
\]
where $T_s$ is the sampling time. Introducing the notation
\[
x(k)=x_c(t_0+kT_s), \qquad u(k)=u_c(t_0+kT_s),
\]
one obtains the discrete-time approximation
\[
x(k+1)=x(k)+T_s g_c(x(k),u(k)), \qquad y(k)=h_c(x(k),u(k)).
\]
Under suitable regularity assumptions and for sufficiently small sampling times, the discrete-time trajectory approximates the continuous-time one well.\\

For linear continuous-time systems,
\[
\dot{x}(t)=A_c x(t)+B_c u(t), \qquad y(t)=C_c x(t)+D_c u(t),
\]
Euler discretization yields
\[
x(k+1)=Ax(k)+Bu(k), \qquad y(k)=Cx(k)+Du(k),
\]
with
\[
A = I + T_s A_c, \qquad B = T_s B_c, \qquad C=C_c, \qquad D=D_c.
\]
This is straightforward and often useful, although it is not exact.

\subsection{Exact discretization of linear systems}

For linear systems under zero-order hold, one can derive an exact discrete-time model. Using the explicit continuous-time solution and assuming that the input remains constant over one sampling interval, one obtains
\[
x(t_{k+1}) = e^{A_c T_s}x(t_k) + \left(\int_{0}^{T_s} e^{A_c(T_s-\tau)}B_c\,d\tau\right) u(t_k).
\]
Hence the exact discrete-time system is
\[
x(k+1)=Ax(k)+Bu(k),
\]
with
\[
A = e^{A_cT_s},
\qquad
B = \int_0^{T_s} e^{A_c(T_s-\tau)}B_c\,d\tau.
\]
If $A_c$ is invertible, then $B$ can also be written as
\[
B = A_c^{-1}(A-I)B_c.
\]
This model is \textbf{exact} in the sense that it predicts the state of the continuous-time system exactly at the sampling instants, provided the zero-order-hold assumption is satisfied.

\subsection{Solution of linear discrete-time systems}

One major advantage of discrete-time linear systems is that their future state can be written explicitly as a function of the initial condition and the input sequence. Starting from
\[
x(k+1)=Ax(k)+Bu(k),
\]
successive substitution gives
\[
x(k+N)=A^N x(k) + \sum_{i=0}^{N-1} A^i B\,u(k+N-1-i).
\]
Thus, for fixed matrices $A$ and $B$, the predicted state over a horizon is a \textbf{linear function of the current state and of the future inputs}. This observation is fundamental in MPC, because it allows the finite-horizon optimization problem to be written directly in terms of the decision variables and the known initial state.\\

The modeling discussion can be summarized as follows. Continuous-time nonlinear models are general and physically natural, but they are difficult to analyze and use directly in predictive control. Linearization gives a local continuous-time approximation around an operating point. Discretization then converts the model into a form that is directly compatible with digital control implementation. For linear systems, exact zero-order-hold discretization is available, while for general nonlinear systems one often relies on approximate methods such as Euler discretization, whose quality depends strongly on the sampling time. These are precisely the model transformations that make standard MPC formulations possible.

\section{Analysis of LTI Discrete-Time Systems}

Once a model has been put into discrete-time LTI form, one can study several structural properties that are central in estimation and control. Among the most important are the behavior of the system under coordinate changes, asymptotic stability, controllability, and observability. These notions form the classical backbone of linear systems theory and play a direct role in MPC.

\subsection{Coordinate transformations}

Consider the discrete-time LTI system
\[
x(k+1)=Ax(k)+Bu(k), \qquad y(k)=Cx(k)+Du(k).
\]
The input-output behavior is completely determined by the initial state and the input sequence. However, the choice of state variables is \textbf{not unique}. Different state representations may generate exactly the same input-output behavior, and some representations may be more convenient than others for analysis or design.\\
If one introduces a new state variable
\[
\tilde{x}(k)=Tx(k),
\]
where $T$ is invertible, then the transformed system becomes
\[
\tilde{x}(k+1)=\tilde{A}\tilde{x}(k)+\tilde{B}u(k), \qquad
y(k)=\tilde{C}\tilde{x}(k)+\tilde{D}u(k),
\]
with
\[
\tilde{A}=TAT^{-1}, \qquad \tilde{B}=TB, \qquad \tilde{C}=CT^{-1}, \qquad \tilde{D}=D.
\]
The input and output remain unchanged; only the internal coordinates have changed. Such transformations are useful because they can simplify the structure of the matrices and thereby reveal important properties more transparently.

\subsection{Stability of linear discrete-time systems}

For the autonomous system
\[
x(k+1)=Ax(k),
\]
\textbf{asymptotic stability} means that every trajectory converges to the origin as $k\to\infty$. A fundamental theorem states that this happens if and only if all eigenvalues of $A$ lie strictly inside the unit circle, that is,
\[
|\lambda_j|<1, \qquad j=1,\dots,n.
\]
This is the discrete-time analogue of the continuous-time condition $\mathrm{Re}(\lambda_j)<0$.\\\

The intuition becomes especially clear when the system can be diagonalized. In suitable coordinates, the system matrix becomes diagonal, and each transformed state component evolves independently as multiplication by the corresponding eigenvalue. If the magnitude of every eigenvalue is less than one, repeated multiplication drives each component to zero. If any eigenvalue has magnitude greater than or equal to one, some component fails to decay. Even when diagonalization is not possible and Jordan forms must be used, the same conclusion remains valid.

\subsection{Controllability and stabilizability}

A system
\[
x(k+1)=Ax(k)+Bu(k)
\]
is \textbf{controllable} if, for any initial state $x(0)$ and any desired target state $x*$, there exists a finite time $N$ and an input sequence $\{u(0), u(1), \dots, u(N-1)\}$ that steers the system exactly to that target, i.e. $x(N)=x*$. In discrete time, this notion is often also called reachability. By expanding the state recursively:
$$x* = x(N) = A^N x(0) + \sum_{i=0}^{N-1} A^i B u(N-1-i) = A^N x(0)+ \begin{bmatrix}
B & AB & \cdots & A^{N-1}B     
\end{bmatrix}
\begin{bmatrix}u(N-1)\\ u(N-2) \\ \vdots \\ u(0)   \end{bmatrix}$$
It follows from the \textbf{Carley-Hamilton} theorem that $A^k$ can be expressed as a linear combination of $I$, $A$, $\dots$, $A^{n-1}$, so that for all $N\ge n$ 
$$\text{range} \begin{bmatrix}
    B & AB & \cdots & A^{N-1}B
\end{bmatrix} = \begin{bmatrix}
    B & AB & \cdots & A^{n-1}B
\end{bmatrix}
$$
We therefore conclude that the reachable states are precisely those that can be generated by the columns of the \textbf{controllability matrix}
\[
\mathcal{C} = \begin{bmatrix} B & AB & \cdots & A^{n-1}B \end{bmatrix},
\]
this means that the reachable states are precisely those in the column space of $\mathcal{C}$. The system is controllable \textbf{if and only if}
\[
\operatorname{rank}(\mathcal{C}) = n.
\]
A weaker notion is \textbf{stabilizability}. A system is stabilizable if the unstable part of the dynamics can still be influenced by the input, even if some already stable modes are uncontrollable. More precisely, the system is stabilizable if all uncontrollable modes are stable. Controllability always implies stabilizability, but the converse need not hold. In many feedback-design settings, stabilizability is the essential requirement.\\
Stabilizability can be checked using the following condition
$$\text{if rank} \biggl(\begin{bmatrix}\lambda_j I - A |B \end{bmatrix}\biggr) = n \quad \forall \lambda_j \in \Lambda_A^+ \Rightarrow (A,B) \text{is stabilizable}$$
where $\Lambda_A^+$ is the set of eigenvalues of $A$ lying on or outside the unit circle.
\subsection{Observability and detectability}

Dual to controllability is the notion of \textbf{observability}. Consider the zero-input system
\[
x(k+1)=Ax(k), \qquad y(k)=Cx(k).
\]
The system is observable if the initial state can be uniquely reconstructed from a finite sequence of output measurements. By stacking the outputs over time, one obtains the linear relation
\[
\begin{bmatrix}
y(0)\\ y(1)\\ \vdots\\ y(n-1)
\end{bmatrix}
=
\begin{bmatrix}
C\\ CA\\ \vdots\\ CA^{n-1}
\end{bmatrix}x(0).
\]
The matrix
\[
\mathcal{O}=
\begin{bmatrix}
C\\ CA\\ \vdots\\ CA^{n-1}
\end{bmatrix}
\]
is called the \textbf{observability matrix}, and the system is observable \textbf{if and only if}
\[
\operatorname{rank}(\mathcal{O})=n.
\]
Again, the use of only the first $n$ blocks is justified by the Cayley--Hamilton theorem.\\

The weaker property corresponding to stabilizability is \textbf{detectability}. A system is detectable if all unobservable modes are stable, so that although the state may not be reconstructed exactly, the estimation error can still be driven to zero asymptotically. Observability implies detectability, but not every detectable system is observable. In practice, detectability is often the minimal property needed to construct a convergent observer or state estimator.\\
Observability can be checked using the following condition
$$\text{if rank} \biggl(\begin{bmatrix}\lambda_j I - A \\ C \end{bmatrix}\biggr) = n \quad \forall \lambda_j \in \Lambda_A^+ \Rightarrow (A,C) \text{ is detectable}$$
where $\Lambda_A^+$ is the set of eigenvalues of $A$ lying on or outside the unit circle.

\begin{theorembox}[Structural properties]
Controllability and observability are exact structural properties of a model. Their relaxed counterparts, stabilizability and detectability, are often sufficient for controller and observer design and are especially relevant when only the unstable modes must be influenced or reconstructed.
\end{theorembox}

\section{Analysis of Nonlinear Discrete-Time Systems}

The previous section dealt with discrete-time LTI systems, for which many results can be stated in terms of matrices and eigenvalues. For nonlinear systems, the situation is more subtle. Stability is no longer merely a property of the matrix $A$, but a \textbf{property of an equilibrium point} of the nonlinear dynamics. The appropriate conceptual framework is Lyapunov stability theory.

\subsection{Lyapunov stability of nonlinear systems}

Consider an autonomous nonlinear discrete-time system
\[
x(k+1)=g(x(k)),
\]
and let $\bar{x}$ be an equilibrium point, meaning that
\[
g(\bar{x})=\bar{x}.
\]
\textbf{Lyapunov stability} formalizes the intuitive idea that if the system is perturbed only slightly from equilibrium, then it should remain close to equilibrium for all future time. More precisely, $\bar{x}$ is Lyapunov stable if for every $\varepsilon>0$ there exists a $\delta(\varepsilon)>0$ such that
\[
\|x(0)-\bar{x}\|<\delta(\varepsilon)
\quad\Longrightarrow\quad
\|x(k)-\bar{x}\|<\varepsilon \qquad \forall k\ge 0.
\]
If, in addition, trajectories converge to the equilibrium ( i.e. is attractive), meaning that
\[
\lim_{k\to\infty}\|x(k)-\bar{x}\|=0
\qquad \forall x(0)\in \R^n,
\]
then the equilibrium is \textbf{globally asymptotically stable}.
\subsection{Lyapunov functions}

A standard way to prove stability is to construct a Lyapunov function. The intuition comes from mechanics: if the total energy of a damped system decreases over time, then the system tends to settle to an equilibrium. Lyapunov functions generalize this idea to arbitrary dynamical systems.\\
For an equilibrium at the origin, a \textbf{Lyapunov function} is a scalar function
\[
V:\R^n\to\R
\]
that is continuous at the origin, finite everywhere, radially unbounded, positive definite, and strictly decreases along system trajectories. In particular, one requires
\[
V(0)=0, \qquad V(x)>0 \ \text{for all } x\neq 0,
\]
\[
\|x\|\to\infty \ \Longrightarrow\ V(x)\to\infty,
\]
and
\[
V(g(x))-V(x)\le -\alpha(x),
\]
where $\alpha$ is a continuous positive definite function. \textbf{If such a function exists, then the origin is globally asymptotically stable}. If the decrease condition holds only with $\alpha$ positive semidefinite, then one obtains Lyapunov stability but not necessarily asymptotic convergence.

\subsection{Indirect Lyapunov method}

Although Lyapunov functions are powerful, they are often \textbf{difficult to construct} for general nonlinear systems. A useful alternative is the indirect Lyapunov method, which studies the linearization around the equilibrium. Let
\[
A = \left.\frac{\partial g}{\partial x^\top}\right|_{x=0},
\]
then the origin is \textbf{locally asymptotically stable if} all eigenvalues of $A$ satisfy
\[
|\lambda_i|<1.
\]
On the other hand, if at least one eigenvalue satisfies
\[
|\lambda_i|>1,
\]
then the origin is unstable. The borderline case in which some eigenvalue lies exactly on the unit circle is \textbf{inconclusive}, and one must resort to a more refined analysis, often through a Lyapunov function.\\

This theorem is especially important because it connects nonlinear local stability with the familiar eigenvalue test from linear systems. It also clarifies why linearization is so useful: it can certify local asymptotic stability when the equilibrium lies safely inside the stable region, but it cannot always settle the question in marginal cases.

\subsection{Lyapunov stability of linear systems}

For \textbf{linear systems}, Lyapunov theory becomes constructive. Consider
\[
x(k+1)=Ax(k).
\]
A natural candidate Lyapunov function is quadratic,
\[
V(x)=x^\top P x,
\]
with $P=P^\top>0$. Evaluating the change along trajectories gives
\[
V(x(k+1))-V(x(k))
= x(k)^\top(A^\top P A - P)x(k).
\]
If one can choose $P>0$ such that
\begin{equation}
A^\top P A - P = -Q,
\qquad Q>0,
\label{eq:lyap_equation}
\end{equation}
then
\[
V(x(k+1))-V(x(k)) = -x(k)^\top Q x(k) < 0
\]
for all nonzero $x(k)$, which proves asymptotic stability. Equation \eqref{eq:lyap_equation} is the \textbf{discrete-time Lyapunov equation}.\\

A central theorem states that, for discrete-time LTI systems, the Lyapunov equation has a \textbf{unique positive definite solution} $P$ for every given $Q>0$ \textbf{if and only if} all eigenvalues of $A$ lie inside the unit circle. Thus, in the linear case, eigenvalue stability and the existence of a quadratic Lyapunov function are equivalent. Moreover, because the system is linear, this stability is automatically global (stability is always "global" for linear systems, since the dynamics are the same everywhere in the state space).\\

There is also an important \textbf{interpretation} of the matrix $P$ in terms of infinite-horizon cost. If the system is asymptotically stable and one defines
\[
\Psi(x(0))=\sum_{k=0}^{\infty} x(k)^\top Q x(k),
\]
then this infinite sum can be written as
\[
\Psi(x(0)) = x(0)^\top P x(0),
\]
where $P$ is precisely the solution of the discrete-time Lyapunov equation.\\
We can proof this by defining $H_k = (A^k)^\top Q (A^k)$ and $P = \sum_{k=0}^\infty H_k$. Notice that the limit of the sum exists because the system is asymptotically stable. Then we have that
$$A^\top H_k A = (A^{k+1})^\top Q (A^{k+1}) = H_{k+1}$$
Thus 
$$ A^\top P A = A^\top \left(\sum_{k=0}^\infty H_k\right) A = \sum_{k=0}^\infty A^\top H_k A = \sum_{k=0}^\infty H_{k+1} = P - H_0 = P- Q$$
In other words, the same quadratic function that certifies stability also represents the infinite-horizon cost-to-go associated with the stage cost $x^\top Qx$. This link is conceptually very important for MPC, because later the optimal cost itself will play the role of a Lyapunov function for the closed-loop system under suitable assumptions.

\section{Summary}

The material of this chapter provides the basic system-theoretic background required for MPC. Starting from nonlinear continuous-time models, one typically linearizes around an operating point and then discretizes the system to obtain a model suitable for digital control implementation. For linear discrete-time systems, explicit predictions over a finite horizon can be written as linear functions of the current state and the future input sequence, which is exactly the structure exploited in predictive control.\\
On the analysis side, stability of LTI systems is characterized by eigenvalue location, while controllability and observability determine whether the system can be influenced and reconstructed from measurements. Their weaker counterparts, stabilizability and detectability, are often sufficient for design tasks. For nonlinear systems, Lyapunov theory provides the appropriate framework to study stability of equilibria. In the linear case, quadratic Lyapunov functions arise naturally through the discrete-time Lyapunov equation and connect stability analysis with infinite-horizon cost. This bridge between dynamics, stability, and cost is one of the conceptual pillars on which the theory of MPC is built.
